\section{Analysis and Discussion}
\label{sec:analysis}

\subsection{Where the aggregate gain comes from}

The specialised parser's \PacsOursCorrectA{} strict successes consist of 549 correct executed
answers and 138 correct explicit abstentions.  In contrast, the untuned model produces only 291
correct answers and 41 correct abstentions; it answers 114 out-of-grammar cases and emits 104
invalid trees.  Typed-program supervision therefore changes both semantic coverage and the
model's willingness to use the unsupported action.  It is not merely a syntax repair: the base
model already emits a high fraction of type-valid trees, yet many encode the wrong role,
projection, or composition.

The cloud planner produces 324 correct answers and 126 correct abstentions, but also 161 invalid
trees, 57 unparseable outputs, and 25 API errors.  Counting service errors as failures is
appropriate for an end-to-end single-run comparison, but it mixes model capability with provider
availability.  A future controlled comparison should repeat the run, report provider failures
separately, and keep the same retry budget for every system.

\subsection{Strict and safe outcomes answer different questions}

The specialised Channel-A run contains 73 literal-faithful empty-result executions.  The strict
metric does not credit them because the benchmark requires an explicit abstention for those
non-answerable controls.  The supplementary safe metric does, raising the score from 74.51\% to
82.43\%.  This gap is informative but should not be presented as a second answer-accuracy number:
it measures whether the model avoided an unsupported factual answer under a broader operational
policy.  The full outcome counts and definition appear in Appendix \ref{app:full-results}.

This distinction also exposes a product decision.  A user asking for a count over an empty scope
might reasonably prefer ``0'', while a missing-operator question requires an explicit statement
that the operation is unsupported.  A production interface should make these policies visible
rather than merge both into a generic refusal label.

\subsection{Composition is not equivalent to depth}

The non-monotonic L1/L2/L3 results show that depth is an incomplete difficulty measure.  An L3
set difference can be structurally regular and well represented in supervision, whereas an L2
grouped concentration query may require a rarer operator and harder return type.  The family
results support reporting a two-dimensional family-by-depth matrix, with cell counts beside
percentages, rather than a single line labelled ``accuracy versus complexity.''  The final
figure brief in the repository specifies this diagnostic and retains the exact table in the
appendix.

Declared unseen shapes also do not imply unseen primitives.  PACS records whether an exact shape
signature occurs in the training export, but it does not calculate atom and compound divergence
as CFQ does.  The 74.95\% unseen-shape score therefore supports generalisation to the benchmark's
declared structural holdout; it does not establish systematic compositional generalisation in
the stronger formal sense.

\subsection{What the full-runtime result can and cannot attribute}

The \RuntimeOurs\% system result is consistent with the intended benefit of specialised planning
inside an execute-and-check runtime.  Category gains on bridge, comparison, and set questions
align with the PACS family improvements.  Nevertheless, the comparison does not isolate the
understander, router, deterministic compiler, graph planner, verifier, or repair loop.  A causal
component claim would require compiler-only, always-planner, conditional-routing, no-verifier,
and no-repair variants on the same snapshot and budget.  We therefore report the observed full
stack difference and avoid assigning it to any single component.

The runtime's strongest methodological property is instead inspectability.  A failure can be
localised to an unparseable intent, routing veto, graph inconsistency, grounding error,
preflight failure, empty intermediate set, postflight defect, or evidence decision.  This is a
practical improvement over direct answer generation even when the final accuracy is unchanged.
It should be called \emph{auditable} rather than \emph{fully interpretable}: the learned
understander and planner do not provide a complete causal explanation of their own decisions.

\subsection{Data quality remains part of model quality}

The bridge example returns 17 distinct stored names, but multiple strings appear to refer to the
same organisational family.  An exact set match can therefore be correct under the benchmark
and still be inconvenient for a user who expects resolved legal entities.  Similarly, a correct
program over an unnormalised currency column can return a reproducible but semantically
mislabelled total.  These are not peripheral data-engineering issues: they define the denotation
against which the model is trained and scored.  Future versions should version the entity and
currency semantics together with the benchmark, rather than updating the KG under an unchanged
test label.
